% PDF 212--216；原书第 184--188 页。
\section{基础——强化学习框架}
\label{sec:ch05-rl}
\begin{keypoint}
\noindent\textbf{为什么用强化学习做大语言模型对齐}

大语言模型擅长生成连贯文本，但并不天然理解细致的人类偏好和社会规范。人能直观地权衡上下文、文化差异和伦理影响等复杂因素，因而擅长评估和选择模型输出。自动化系统往往难以充分理解并顾及这些因素。在这种情况下，以传统损失函数和常规优化技术训练语言模型，使其符合人类偏好，并不实际。原因是人的反馈分数就像图~\ref{fig:ch05-feedback} 所示的黑箱。通常的神经网络会对分数反向传播，但在这里无法这样做，因为那需要计算作出主观文本评估的人类反馈机制的梯度。强化学习（reinforcement learning，RL）是可处理不可微学习信号的技术之一，已成为在大语言模型调优中纳入人类偏好的常用技术。
\end{keypoint}

\begin{figure}[!htbp]
\centering
% 重绘自原书第 185 页的图 5.1。
\begin{tikzpicture}[
  font=\footnotesize\sffamily,
  flow/.style={draw=black!68,-{Stealth[length=1.35mm]},line width=.52pt},
  layer/.style={draw=black!64,rounded corners=2pt,align=center,minimum width=31mm,
    minimum height=5.0mm,inner sep=1.4pt,fill=green!15,font=\scriptsize\sffamily}]
\node[layer,fill=yellow!35] (softmax) at (2.2,5.02) {Softmax};
\node[layer,fill=yellow!35] (linear) at (2.2,4.48) {线性层};
\node[layer] (norm3) at (2.2,3.90) {相加与归一化};
\node[layer,fill=green!24] (ffn) at (2.2,3.32) {逐位置前馈网络};
\node[layer] (norm2) at (2.2,2.74) {相加与归一化};
\node[layer,fill=green!33] (cross) at (2.2,2.16) {编码器—解码器注意力};
\node[layer] (norm1) at (2.2,1.58) {相加与归一化};
\node[layer,fill=green!40] (mask) at (2.2,.99) {掩码多头注意力};
\node[draw=black!60,densely dashed,rounded corners=4pt,fit=(norm3)(mask),
      inner sep=3.6pt] (decoder) {};
\node[layer,fill=cyan!13] (embed) at (2.2,-.17) {嵌入};
\node[draw,circle,inner sep=0pt,minimum size=3.8mm] (plus) at (2.2,.39) {$+$};
\node[layer,fill=blue!25,minimum width=16mm] (pos) at (4.72,.39) {位置编码};
\node[anchor=west,align=left] (lab) at (4.02,4.66) {大语言模型\\智能体};
\node[anchor=west,font=\scriptsize\sffamily,align=left] at (4.01,2.17) {$\times n$ 层\\解码器};
\begin{scope}[on background layer]
\node[draw=blue!30!black!62,fill=blue!8,rounded corners=9pt,
      fit=(softmax)(embed)(pos)(lab),inner sep=5pt] (model) {};
\end{scope}
\foreach \first/\second in {embed/plus,plus/mask,mask/norm1,norm1/cross,cross/norm2,
     norm2/ffn,ffn/norm3,norm3/linear,linear/softmax}{
  \draw[flow] (\first) -- (\second);
}
\draw[flow] (pos.west) -- (plus.east);
\foreach \first/\second in {mask/norm1,cross/norm2,ffn/norm3}{
  \draw[flow] (\first.south east) -- ++(.20,-.08) |- (\second.east);
}
\node[align=center] at (-4.19,3.89) {人工标注者};
\fill[orange!63!black!70] (-4.19,3.16) circle[radius=2.8mm];
\fill[orange!63!black!70] (-4.79,2.31) -- (-4.79,2.56)
  .. controls (-4.74,3.02) and (-3.64,3.02) .. (-3.59,2.56)
  -- (-3.59,2.31) -- cycle;
\foreach \x/\mark/\shade in {-4.54/\checkmark/green,-3.84/\times/red}{
  \draw[blue!50!black] (\x-.19,1.89) rectangle (\x+.19,1.25);
  \foreach \y in {1.42,1.54,1.66,1.78}{
    \draw[blue!50!black,line width=.32pt] (\x-.12,\y) -- (\x+.12,\y);
  }
  \node[\shade!61!black] at (\x,1.01) {$\mark$};
}
\draw[flow] (-3.20,3.35) to[bend left=29] node[below] {反馈} (.04,3.35);
\draw[flow] (.04,1.61) to[bend left=29] node[above] {模型输出} (-3.20,1.61);
\node[draw=blue!50!black,fill=black!9,align=left,font=\scriptsize\sffamily,
  inner sep=2.2pt] at (-1.55,4.46) {\textcolor{red}{$\times$}\enspace 不可微，无梯度};
\end{tikzpicture}

\caption{人工评估模型输出虽然有效，但反馈不可微，无法用传统的基于梯度的方法直接训练大语言模型。}
\label{fig:ch05-feedback}
\end{figure}

强化学习的发展可追溯到多个学科的方法和理论。心理学家 B. F. Skinner 的研究是其中的早期来源。他提出操作性条件作用（operant conditioning）的概念，强调奖励与惩罚如何塑造生物的行为（\refcite{ch05-skinner}{Skinner, 1965}）。这一思想为在计算模型中研究试错学习奠定了基础。

本节将讨论强化学习的基本概念，说明其含义，并给出相应的数学形式和公式。为帮助理解，采用图~\ref{fig:ch05-maze} 所示的简单迷宫，说明智能体如何运用强化学习原理学会在迷宫中移动并到达目标。设智能体在一个由方格组成的迷宫里移动。迷宫有起点和终点，墙或隔板形成各种障碍。智能体的目标是避开障碍，找到从起点到终点最短、最高效的路径。

\begin{figure}[!htbp]
\centering
% 重绘自原书第 186 页的图 5.2。
\begin{tikzpicture}[
  x=.61cm,y=.61cm,font=\footnotesize\sffamily,
  wall/.style={draw=black,line width=3.25pt,line cap=round,line join=round},
  guide/.style={draw=red!48,-{Stealth[length=1.3mm]},line width=.55pt},
  callout/.style={draw=black!18,rounded corners=2pt,fill=black!5,align=center,inner sep=3pt}]
\draw[wall] (1,5)--(5,5)--(5,0)--(4,0);
\draw[wall] (3,0)--(0,0)--(0,5);
\draw[wall] (2,5)--(2,4) (4,5)--(4,4) (0,3)--(3,3)--(3,4) (4,3)--(5,3);
\draw[wall] (1,0)--(1,2)--(3,2)--(3,1) (2,1)--(4,1)--(4,2);
\foreach \x/\y/\a in {.50/5.32/0,4.50/4.45/0,3.50/2.48/-90,3.50/-.34/0}{
  \begin{scope}[shift={(\x,\y)},rotate=\a,scale=.67]
  \draw[fill=black!12,line width=.35pt] (-.20,-.34) rectangle (.20,-.05);
  \draw[fill=black!13,line width=.35pt,rounded corners=.8pt] (-.22,-.02) rectangle (.22,.29);
  \draw[line width=.35pt] (0,.29)--(0,.40) (-.20,-.13)--(-.34,-.22) (.20,-.13)--(.34,-.22);
  \fill (-.08,.16) circle[radius=.45pt] (.08,.16) circle[radius=.45pt] (0,.40) circle[radius=.5pt];
  \draw[line width=.35pt] (-.11,-.34)--(-.11,-.43) (.11,-.34)--(.11,-.43);
  \end{scope}
}
\node[callout,fill=red!15] (start) at (-2.24,5.32) {智能体起点};
\node[above=1.5pt of start] {状态};
\node[callout] (open) at (-2.72,2.48) {智能体面向空地};
\node[above=1.5pt of open] {状态};
\node[callout] (closed) at (7.67,4.45) {智能体面向墙};
\node[above=1.5pt of closed] {状态};
\node[callout,fill=green!18] (end) at (7.42,-.34) {智能体终点};
\node[above=1.5pt of end] {状态};
\draw[guide] (start.east)--(.19,5.32);
\draw[guide] (open.east)--(3.16,2.48);
\draw[guide] (closed.west)--(4.80,4.45);
\draw[guide] (end.west)--(3.82,-.34);
\end{tikzpicture}

\caption{强化学习提供一种机制，可以奖励使智能体更接近迷宫出口的决策。}
\label{fig:ch05-maze}
\end{figure}

每一步，智能体都收到一个状态 $s$。状态可以包含正前方是墙还是空地，左右两侧是墙还是空地，在邻近通道中能看到多远才遇到墙，以及智能体到达这里之前所做的移动。对每个这样的状态，智能体都可以从有限的动作集 $A$ 中选择动作，例如向上、向下、向左或向右移动。根据采取的动作 $a$，智能体获得奖励或惩罚 $r$，据此学习。例如，到达终点可获得正奖励。为鼓励效率，每走一步可给一个较小的负奖励。撞墙或走出迷宫边界时，则给一个绝对值较大的负奖励。

起初，智能体不知道迷宫的布局，也不知道最佳路径。探索环境时，它遇到表示迷宫中不同位置的状态，采取各种动作，再进入新的状态。对这些动作反复施加奖励或惩罚，会影响智能体今后在各状态下给每种可能动作分配的概率。如果强化学习成功，学到的这些概率会使智能体比起初更高效地走出迷宫。

马尔可夫决策过程（Markov decision process，MDP）是强化学习的基础数学框架，用离散时间随机控制过程为情境建模（\refcite{ch05-puterman}{Puterman, 1990}）。

如图~\ref{fig:ch05-mdp} 所示，在马尔可夫决策过程中，作决策的智能体与周围环境按时间顺序反复交互。在每个离散时刻，智能体获取环境状态的一种表示，并据此选择动作。随后，环境转移到新状态，智能体根据先前动作造成的结果获得奖励。智能体在这个过程中主要希望最大化在不同状态下采取动作而得到的累计奖励。

\begin{figure}[!htbp]
\centering
% 重绘自原书第 187 页的图 5.3。
\begin{tikzpicture}[
  font=\footnotesize\sffamily,
  nodebox/.style={draw=blue!42!black,fill=cyan!16,rounded corners=3pt,align=center,
    minimum width=23mm,minimum height=9mm,inner sep=2pt},
  pathline/.style={draw=black!68,-{Stealth[length=1.35mm]},line width=.55pt}]
\node[nodebox] (agent) at (0,0) {智能体};
\node[nodebox,fill=yellow!24] (env) at (0,-2.78) {环境};
\draw[pathline] (env.west) -- (-3.12,-2.78) -- (-3.12,0) -- (agent.west);
\draw[pathline,densely dashed] ([yshift=1.7pt]env.west) --
  (-2.24,-2.72) -- (-2.24,.13) -- ([yshift=3.7pt]agent.west);
\draw[pathline,densely dotted] (agent.east) -- (3.07,0) -- (3.07,-2.78) -- (env.east);
\node[anchor=east] at (-3.23,-1.40) {1.\ 状态 $S_t$};
\node[anchor=west] at (3.13,-1.40) {2.\ 动作 $A_t$};
\node[anchor=west] at (-2.10,-1.40) {3.\ 状态 $S_{t+1}$，奖励 $R_{t+1}$};
\end{tikzpicture}

\caption{强化学习的马尔可夫决策过程。}
\label{fig:ch05-mdp}
\end{figure}

理解这一过程需要几个术语。

\begin{itemize}
\item \textbf{状态（$S$）和动作（$A_t$）。}在 MDP 中，状态表示系统的配置，动作是决策者在各状态下可作的选择。时刻 $t$ 的状态和所采取的动作分别记为 $S_t$ 和 $A_t$。
\item \textbf{奖励（$R_t$）。}奖励是某一动作引起状态转移时得到的实数。奖励函数记作 $R_t$，量化在给定状态下选择某一动作所带来的即时收益。
\item \textbf{转移（$P$）。}转移函数 $P(S_{t+1}\mid S_t,A_t)$ 描述给定当前状态和动作时下一状态的概率分布。它刻画环境如何随智能体的动作而变化。
\item \textbf{策略（$\pi$）。}策略 $\pi$ 定义为将状态映射到动作概率分布的函数。形式上，$\pi(A_t\mid S_t)$ 决定了智能体在状态 $S_t$ 下采取各个动作的可能性。
\end{itemize}

智能体与环境交互，产生一系列动作和状态，所遵循的策略影响这一序列。状态与动作的轨迹描述了这个迭代过程。智能体按策略 $\pi$ 执行动作。要使未来奖励最大，策略需要描述在每个状态下怎样选择最优动作。

\begin{itemize}
\item \textbf{轨迹。}轨迹是智能体经历的状态和动作序列 $\{S_0,A_0,\allowbreak S_1,A_1,\allowbreak\ldots,\allowbreak S_T,A_T\}$，通常止于终止状态，标志着一个回合的结束。
\item \textbf{回报（$G_t$）。}回报是使用折扣因子计算的一条轨迹上的累计奖励。
\[
G_t=\sum_{k=0}^{T}\gamma^k R_{t+k}.
\]
其中，$\gamma$ 是折扣因子，降低较晚获得的奖励的价值，使即时奖励具有较大的相对价值。折扣因子 $\gamma$ 的取值介于 0 与 1 之间，在回报计算中起重要作用，因为它会折减未来奖励的价值，体现对即时奖励相对于延迟奖励的偏好。
\end{itemize}

关于强化学习、其组成部分及算法的详细讨论，见附录~\ref{app:reinforcement-learning}。

\section{将强化学习框架用于基于人类反馈的大语言模型调优}
\label{sec:ch05-map-rlhf}
下面说明强化学习框架中的状态、动作、策略、奖励模型等要素，如何对应大语言模型用人类反馈做对齐调优的过程。

\begin{enumerate}
\item \textbf{智能体。}智能体指语言模型本身。它与环境交互，根据输入状态采取动作，并根据收到的反馈，即奖励，进行学习。
\item \textbf{状态。}状态是提供给模型的上下文，通常以输入提示的形式给出。例如，若输入为“\texttt{ChatGPT is one of the large languages}”，这段文字定义了当前状态。
\item \textbf{动作。}动作是模型根据当前状态选择的下一个词元或词。例如，给定上面的提示，模型可能预测几个候选词，如“model”“tools”或“systems”。选择其中一个接到末尾，就是所采取的动作。
\item \textbf{奖励模型。}语言模型按输出质量获得奖励。准确、切题、有帮助、无害且连贯的“好回答”得到奖励。不准确、不切题、有害或不连贯的“坏回答”得到零奖励或负奖励。
\item \textbf{策略。}在语言模型的语境下，策略实质上就是语言模型本身。语言模型通过为给定状态（输入提示）下各可能动作（后续词元）的概率分布建模，定义了策略。
\end{enumerate}
